\documentclass[11pt,a4paper]{article}
% Preambulo comun de los reportes de Tecnologias Celulares.
% Cada main.tex hace: \documentclass ... % Preambulo comun de los reportes de Tecnologias Celulares.
% Cada main.tex hace: \documentclass ... % Preambulo comun de los reportes de Tecnologias Celulares.
% Cada main.tex hace: \documentclass ... \input{../comun/preambulo.tex}
\usepackage[utf8]{inputenc}
\usepackage[T1]{fontenc}
\usepackage[spanish,mexico]{babel}
\usepackage{graphicx}
\usepackage[margin=2.5cm]{geometry}
\usepackage{amsmath}
\usepackage{amssymb}
\usepackage{booktabs}
\usepackage{tabularx}
\usepackage{array}
\usepackage{float}
\usepackage{setspace}
\usepackage{parskip}
\usepackage{fancyhdr}
\usepackage{textcomp}
\usepackage{xcolor}
\usepackage{caption}
\usepackage{fvextra}
\usepackage[hidelinks]{hyperref}

\onehalfspacing
\setlength{\parskip}{6pt}
\sloppy
\setlength{\emergencystretch}{3em}
\setlength{\headheight}{40pt}
\captionsetup{font=small,labelfont=bf}
\setlength{\tabcolsep}{6pt}

\pagestyle{fancy}
\fancyhf{}
\fancyhead[L]{\includegraphics[width=1cm]{/home/yovick/repos/UAZ-IEI/Assets/UIE.png}}
\fancyhead[R]{\small Tecnologías Celulares}
\fancyfoot[C]{\thepage}
\renewcommand{\headrulewidth}{0.4pt}

\DefineVerbatimEnvironment{codigo}{Verbatim}{%
  breaklines=true,breakanywhere=true,fontsize=\small,frame=single,
  framesep=3pt,rulecolor=\color{gray!50}}

% Figura centrada con pie: \figura[ancho]{archivo}{leyenda}
\newcommand{\figura}[3][0.85]{%
  \begin{center}
    \includegraphics[width=#1\linewidth]{#2}\\[2pt]
    {\small #3}
  \end{center}}

% Portada: \portada{numero}{titulo de la practica}
\newcommand{\portada}[2]{%
\begin{titlepage}
  \thispagestyle{empty}
  \begin{center}
    \vspace*{1cm}
    \includegraphics[width=4.5cm]{/home/yovick/repos/UAZ-IEI/Assets/uaz.png}
    \vspace{0.8cm}

    {\large Universidad Autónoma de Zacatecas}\\[4pt]
    {\large Unidad Académica de Ingeniería Eléctrica}\\[4pt]
    {\normalsize Ingeniería Electrónica Industrial con especialidad en Telecomunicaciones}

    \vspace{1.8cm}
    {\LARGE\bfseries Práctica #1}\\[10pt]
    {\Large\bfseries #2}
  \end{center}

  \vspace{2.2cm}
  \begin{center}
    \begin{tabular}{rl}
      \textbf{Materia:} & Tecnologías Celulares\\
      \textbf{Docente:} & Dr. Víktor Iván Rodríguez Abdalá\\
      \textbf{Alumno:}  & Arnold Yovick Ríos Zaldívar\\
      \textbf{Fecha:}   & Noviembre de 2026\\
    \end{tabular}
  \end{center}
  \vfill
\end{titlepage}}

\usepackage[utf8]{inputenc}
\usepackage[T1]{fontenc}
\usepackage[spanish,mexico]{babel}
\usepackage{graphicx}
\usepackage[margin=2.5cm]{geometry}
\usepackage{amsmath}
\usepackage{amssymb}
\usepackage{booktabs}
\usepackage{tabularx}
\usepackage{array}
\usepackage{float}
\usepackage{setspace}
\usepackage{parskip}
\usepackage{fancyhdr}
\usepackage{textcomp}
\usepackage{xcolor}
\usepackage{caption}
\usepackage{fvextra}
\usepackage[hidelinks]{hyperref}

\onehalfspacing
\setlength{\parskip}{6pt}
\sloppy
\setlength{\emergencystretch}{3em}
\setlength{\headheight}{40pt}
\captionsetup{font=small,labelfont=bf}
\setlength{\tabcolsep}{6pt}

\pagestyle{fancy}
\fancyhf{}
\fancyhead[L]{\includegraphics[width=1cm]{/home/yovick/repos/UAZ-IEI/Assets/UIE.png}}
\fancyhead[R]{\small Tecnologías Celulares}
\fancyfoot[C]{\thepage}
\renewcommand{\headrulewidth}{0.4pt}

\DefineVerbatimEnvironment{codigo}{Verbatim}{%
  breaklines=true,breakanywhere=true,fontsize=\small,frame=single,
  framesep=3pt,rulecolor=\color{gray!50}}

% Figura centrada con pie: \figura[ancho]{archivo}{leyenda}
\newcommand{\figura}[3][0.85]{%
  \begin{center}
    \includegraphics[width=#1\linewidth]{#2}\\[2pt]
    {\small #3}
  \end{center}}

% Portada: \portada{numero}{titulo de la practica}
\newcommand{\portada}[2]{%
\begin{titlepage}
  \thispagestyle{empty}
  \begin{center}
    \vspace*{1cm}
    \includegraphics[width=4.5cm]{/home/yovick/repos/UAZ-IEI/Assets/uaz.png}
    \vspace{0.8cm}

    {\large Universidad Autónoma de Zacatecas}\\[4pt]
    {\large Unidad Académica de Ingeniería Eléctrica}\\[4pt]
    {\normalsize Ingeniería Electrónica Industrial con especialidad en Telecomunicaciones}

    \vspace{1.8cm}
    {\LARGE\bfseries Práctica #1}\\[10pt]
    {\Large\bfseries #2}
  \end{center}

  \vspace{2.2cm}
  \begin{center}
    \begin{tabular}{rl}
      \textbf{Materia:} & Tecnologías Celulares\\
      \textbf{Docente:} & Dr. Víktor Iván Rodríguez Abdalá\\
      \textbf{Alumno:}  & Arnold Yovick Ríos Zaldívar\\
      \textbf{Fecha:}   & Noviembre de 2026\\
    \end{tabular}
  \end{center}
  \vfill
\end{titlepage}}

\usepackage[utf8]{inputenc}
\usepackage[T1]{fontenc}
\usepackage[spanish,mexico]{babel}
\usepackage{graphicx}
\usepackage[margin=2.5cm]{geometry}
\usepackage{amsmath}
\usepackage{amssymb}
\usepackage{booktabs}
\usepackage{tabularx}
\usepackage{array}
\usepackage{float}
\usepackage{setspace}
\usepackage{parskip}
\usepackage{fancyhdr}
\usepackage{textcomp}
\usepackage{xcolor}
\usepackage{caption}
\usepackage{fvextra}
\usepackage[hidelinks]{hyperref}

\onehalfspacing
\setlength{\parskip}{6pt}
\sloppy
\setlength{\emergencystretch}{3em}
\setlength{\headheight}{40pt}
\captionsetup{font=small,labelfont=bf}
\setlength{\tabcolsep}{6pt}

\pagestyle{fancy}
\fancyhf{}
\fancyhead[L]{\includegraphics[width=1cm]{/home/yovick/repos/UAZ-IEI/Assets/UIE.png}}
\fancyhead[R]{\small Tecnologías Celulares}
\fancyfoot[C]{\thepage}
\renewcommand{\headrulewidth}{0.4pt}

\DefineVerbatimEnvironment{codigo}{Verbatim}{%
  breaklines=true,breakanywhere=true,fontsize=\small,frame=single,
  framesep=3pt,rulecolor=\color{gray!50}}

% Figura centrada con pie: \figura[ancho]{archivo}{leyenda}
\newcommand{\figura}[3][0.85]{%
  \begin{center}
    \includegraphics[width=#1\linewidth]{#2}\\[2pt]
    {\small #3}
  \end{center}}

% Portada: \portada{numero}{titulo de la practica}
\newcommand{\portada}[2]{%
\begin{titlepage}
  \thispagestyle{empty}
  \begin{center}
    \vspace*{1cm}
    \includegraphics[width=4.5cm]{/home/yovick/repos/UAZ-IEI/Assets/uaz.png}
    \vspace{0.8cm}

    {\large Universidad Autónoma de Zacatecas}\\[4pt]
    {\large Unidad Académica de Ingeniería Eléctrica}\\[4pt]
    {\normalsize Ingeniería Electrónica Industrial con especialidad en Telecomunicaciones}

    \vspace{1.8cm}
    {\LARGE\bfseries Práctica #1}\\[10pt]
    {\Large\bfseries #2}
  \end{center}

  \vspace{2.2cm}
  \begin{center}
    \begin{tabular}{rl}
      \textbf{Materia:} & Tecnologías Celulares\\
      \textbf{Docente:} & Dr. Víktor Iván Rodríguez Abdalá\\
      \textbf{Alumno:}  & Arnold Yovick Ríos Zaldívar\\
      \textbf{Fecha:}   & Noviembre de 2026\\
    \end{tabular}
  \end{center}
  \vfill
\end{titlepage}}


\begin{document}

\portada{7}{Uso del prefijo cíclico para eliminar la dispersión del canal}

\textbf{Script:} \texttt{practica7.py} (Python con NumPy)

\section{Objetivo}

Analizar el efecto de la dispersión temporal de un canal multitrayectoria sobre señales
OFDM, implementar el prefijo cíclico y comparar el desempeño del sistema con y sin prefijo
mediante la tasa de error de bit.

\section{Fundamento teórico}

Un canal multitrayectoria se modela con una respuesta impulsional finita
\begin{equation}
  h[n] = \sum_{l=0}^{L-1} h_l\,\delta[n-l].
\end{equation}
La señal recibida es la convolución lineal $y[n]=h[n]*x[n]+n[n]$. En OFDM, la convolución
lineal entre símbolos consecutivos produce interferencia entre símbolos (ISI).

El prefijo cíclico copia las últimas $N_{cp}$ muestras del símbolo OFDM al inicio del
mismo. Si $N_{cp}\ge L-1$, la parte útil del bloque recibido contiene una convolución
\emph{circular} con $h[n]$, y entonces
\begin{equation}
  Y[k] = H[k]\,X[k] + N[k],
  \qquad
  H[k] = \sum_{l=0}^{L-1} h_l\,e^{-j2\pi kl/N}.
\end{equation}
Cada subportadora queda multiplicada por un solo coeficiente complejo y se puede ecualizar
con una división:
\begin{equation}
  \hat{X}[k] = \frac{Y[k]}{H[k]}.
\end{equation}
El precio es una pérdida de eficiencia espectral de $N_{cp}/(N+N_{cp})$ y una cierta
amplificación del ruido en las subportadoras donde $|H[k]|$ es pequeño.

\section{Desarrollo experimental}

\subsection{Canal utilizado}

El manual propone $h=[1,\,0.6,\,0.3]$. Con tres taps adyacentes la ISI afecta solo 2 de
las 64 muestras de cada bloque (3\,\%), y el efecto del prefijo resulta marginal. Para que
el experimento muestre el fenómeno se conserva la misma energía pero se reparten los ecos
en las muestras 0, 5 y 12:
\begin{equation}
  h = [1,\ 0,\ 0,\ 0,\ 0,\ 0.6,\ 0,\ 0,\ 0,\ 0,\ 0,\ 0,\ 0.3],
\end{equation}
con un retardo máximo de 12 muestras, menor que $N_{cp}=16$. Se reporta además el caso del
manual como referencia.

\subsection{Parámetros}

\begin{tabularx}{\linewidth}{@{}*{2}{>{\raggedright\arraybackslash}X}@{}}
  \toprule
  Parámetro & Valor \\
  \midrule
  $N_{\mathrm{FFT}}$ & 64 \\
  Prefijo cíclico & $N_{cp}=16$ \\
  Modulación & QPSK \\
  Bits por trama & 4096 \\
  Tramas por punto & 30 (122\,880 bits) \\
  Ecualización & zero-forcing por subportadora \\
  \bottomrule
\end{tabularx}

\subsection{Procedimiento}

\begin{enumerate}
  \item Se genera la señal OFDM y se transmite por el canal con \texttt{np.convolve}.
  \item Caso sin prefijo: se toman bloques de 64 muestras, se aplica la FFT y se ecualiza
        con $H[k]$.
  \item Caso con prefijo: se añaden 16 muestras al inicio de cada símbolo, se transmite,
        se eliminan las 16 muestras del prefijo y se aplica FFT y ecualización.
  \item Se calcula la BER para $E_b/N_0$ de 0 a 20 dB y se grafican las constelaciones a
        15 dB.
\end{enumerate}

\section{Código}

\begin{codigo}
def transmitir_con_cp(x, h, ncp=16):
    tx = np.vstack((x[-ncp:], x)).flatten(order="F")
    return tx, np.convolve(tx, h)[: tx.size]

def recibir_con_cp(rx, h, ncp=16):
    bloques = rx.reshape(-1, NFFT + ncp).T
    y = np.fft.fft(bloques[ncp:], axis=0) / np.sqrt(NFFT)
    hh = np.fft.fft(h, NFFT) / np.sqrt(NFFT)
    return (y / hh[:, None]).flatten(order="F")
\end{codigo}

\section{Resultados}

\subsection{Canal}

\figura{figs/canal.png}{Canal del manual y canal dispersivo utilizado en la comparación.}

\subsection{Interferencia entre símbolos}

\figura{figs/tiempo_isi.png}{Señal recibida sin prefijo (arriba) y con prefijo (abajo). Las zonas sombreadas marcan las muestras invadidas por el símbolo anterior.}

Sin prefijo, las 12 primeras muestras de cada bloque contienen la cola del símbolo
anterior; con prefijo, esa cola cae dentro de las 16 muestras descartadas. La discontinuidad
en los bordes de bloque es la ISI que la FFT ya no puede separar.

\subsection{Constelaciones}

\figura{figs/constelaciones.png}{Constelaciones recibidas a $E_b/N_0=15$ dB, sin y con prefijo cíclico.}

Sin prefijo la constelación aparece dispersa y con puntos alejados de las cuatro esquinas;
con prefijo los símbolos se concentran alrededor de $(\pm1,\pm1)/\sqrt{2}$, aunque con
cierta amplificación de ruido en las subportadoras de baja ganancia.

\subsection{BER}

\begin{tabular}{@{}ccc@{}}
  \toprule
  $E_b/N_0$ (dB) & BER sin prefijo & BER con prefijo \\
  \midrule
  0  & 0.10460 & 0.09242 \\
  4  & 0.05374 & 0.04174 \\
  8  & 0.03005 & 0.01755 \\
  12 & 0.02270 & 0.00802 \\
  16 & 0.01937 & 0.00339 \\
  20 & 0.01830 & 0.00068 \\
  \bottomrule
\end{tabular}

\figura{figs/ber.png}{Curvas de BER con y sin prefijo cíclico.}

Sin prefijo la BER se estanca alrededor de $1.8\times10^{-2}$: es un piso de ISI que no
disminuye al aumentar la potencia. Con prefijo la BER sigue bajando hasta
$6.8\times10^{-4}$ a 20 dB. A 14 dB la reducción es de 0.0201 a 0.0054 (factor 3.7) y a
20 dB de 0.0183 a 0.00068 (factor 27).

\subsection{Casos de referencia}

\begin{tabularx}{\linewidth}{@{}*{2}{>{\raggedright\arraybackslash}X}@{}}
  \toprule
  Caso & BER \\
  \midrule
  Canal del manual $[1,\,0.6,\,0.3]$, sin y con prefijo, 15 dB & 0 en ambos \\
  Canal dispersivo sin prefijo y sin ecualizar, 15 dB & 0.0718 \\
  Prefijo insuficiente: eco en la muestra 20 $>N_{cp}=16$, 20 dB & 0.0025 \\
  \bottomrule
\end{tabularx}

Con el canal de tres taps adyacentes el prefijo no cambia la BER medible, porque la ISI
solo afecta 2 muestras por bloque. Si además se omite la ecualización, la BER sube a
0.0718 por la rotación de fase de $H[k]$, no por la ISI. Cuando el eco llega después del
prefijo, la BER vuelve a subir a 0.0025 y aparece un piso: el prefijo solo es efectivo si
cubre el retardo máximo del canal.

\section{Análisis de resultados}

\begin{itemize}
  \item \textbf{¿Qué efecto produce la multitrayectoria?} Cada eco llega con un retardo y
        una atenuación distintos, de modo que la señal recibida es la suma de varias
        réplicas. En el dominio del tiempo eso alarga la respuesta y hace que un símbolo
        invada al siguiente; en frecuencia produce desvanecimientos selectivos.
  \item \textbf{¿Cómo afecta la ISI a los símbolos OFDM?} Rompe la ortogonalidad: la FFT
        del bloque ya no corresponde a la convolución circular, así que la energía del
        símbolo anterior se reparte como error en todas las subportadoras. El resultado es
        un piso de BER que no se corrige aumentando la potencia.
  \item \textbf{¿Por qué el prefijo cíclico mejora el desempeño?} Porque absorbe la cola
        del símbolo anterior dentro de las muestras que el receptor descarta. La parte útil
        queda como una convolución circular, que en frecuencia es una multiplicación por
        $H[k]$ y se puede ecualizar con una sola división.
  \item \textbf{¿Qué ocurre si el retardo del canal supera la longitud del prefijo?} La
        cola invade las muestras útiles y reaparece la ISI. En la simulación, un eco en la
        muestra 20 con $N_{cp}=16$ elevó la BER a 0.0025 a 20 dB, frente a 0.00068 con el
        eco dentro del prefijo.
\end{itemize}

\section{Cuestionario}

\begin{enumerate}
  \item \textbf{¿Qué es la interferencia entre símbolos?} La contaminación de un símbolo
        por la energía de los símbolos anteriores o posteriores, causada por la memoria
        (dispersión temporal) del canal.
  \item \textbf{¿Cuál es la función del prefijo cíclico?} Extender cada símbolo con una
        copia de su cola para que, tras la convolución con el canal, la parte útil se
        comporte como una convolución circular y se pueda ecualizar subportadora por
        subportadora.
  \item \textbf{¿Por qué LTE utiliza prefijo cíclico?} Porque el canal móvil es
        multitrayectoria y el prefijo permite ecualizar en frecuencia con bajo costo,
        manteniendo la ortogonalidad entre subportadoras. LTE usa prefijo normal
        ($\approx4.7\,\mu$s) y extendido ($\approx16.7\,\mu$s) según el entorno.
  \item \textbf{¿Qué relación existe entre la longitud del prefijo y el retardo máximo del
        canal?} Debe cumplirse $N_{cp}\ge L-1$ (en muestras), es decir, el prefijo debe
        durar al menos el retardo máximo del canal. Si no, la ISI reaparece.
  \item \textbf{¿Qué costo tiene agregar un prefijo cíclico?} Reduce la eficiencia
        espectral en $N_{cp}/(N+N_{cp})$ y consume potencia de transmisión en muestras que
        el receptor descarta. También puede amplificar el ruido en las subportadoras con
        desvanecimiento profundo al ecualizar.
\end{enumerate}

\section{Conclusiones}

Se comparó un sistema OFDM con y sin prefijo cíclico sobre un canal multitrayectoria con
retardos en 0, 5 y 12 muestras. Sin prefijo, la ISI produjo un piso de BER de
$1.8\times10^{-2}$ que no bajó al aumentar la potencia. Con un prefijo de 16 muestras y
ecualización zero-forcing, la BER descendió hasta $6.8\times10^{-4}$ a 20 dB. El
experimento también mostró que el prefijo deja de ser efectivo cuando el eco supera su
longitud, y que la ecualización por subportadora es indispensable para explotar la
ortogonalidad.

\end{document}
